\chapter{Special-purpose components}
\index{Library!Components!misc}

The chapter deals with components that are not easily included
in any of the other chapters because of their special nature,
but which are still part of the \MCS system.

One part of these components deals with splitting simulations
into two (or more) stages. For example, a guide system is often
not changed much, and a long simulation of neutron rays
``surviving'' through the guide system could be reused
for several simulations of the instrument back-end, speeding up
the simulations by (typically) one or two orders of magnitude.
As of \MCS 3.x, the recommended components for this are
\textbf{MCPL\_input} and \textbf{MCPL\_output}, which store and read back
neutron rays using the portable, compressed, multi-code MCPL event format
(with backends for \MCS/McXtrace as well as several other Monte Carlo
transport codes -- MCNP, PHITS, Geant4, Vitess, SIMRES). The older,
code-specific event components (\textbf{Virtual\_input}/\textbf{Virtual\_output}
and similar MCNP/Tripoli4/Vitess-specific readers/writers) perform the same
role but are now in the \textbf{obsolete} component category (see the
Obsolete chapter) and are kept only for backward compatibility; new
instrument work should prefer MCPL.

\textbf{Progress\_bar} is a simulation utility that displays the simulation
status, but assumes the form of a component.

Components for simulating instrument resolution functions
(\textbf{Res\_sample}, \textbf{TOF\_Res\_sample} and \textbf{Res\_monitor})
are documented in the Samples and Monitors chapters respectively, matching
their current component-library category (they predate the present
category split and were historically documented here).

\newpage
\section{MCPL: splitting simulations via portable neutron event files}
\label{s:mcpl}
\section{The \texttt{MCPL\_input} McStas Component}
Source-like component that reads neutron state parameters from an mcpl-file.

\subsection*{Identification}
\begin{itemize}
  \item \textbf{Site:} 
  \item \textbf{Author:} Erik B Knudsen
  \item \textbf{Origin:} DTU Physics
  \item \textbf{Date:} Mar 2016
\end{itemize}

\subsection*{Description}
\begin{lstlisting}
Source-like component that reads neutron state parameters from a binary mcpl-file.

MCPL is short for Monte Carlo Particle List, and is a new format for sharing events
between e.g. MCNP(X), Geant4 and McStas.

When used with MPI, the --ncount given on the commandline is overwritten by
#MPI nodes x #events in the file.

Example: MCPL_input(filename=voutput,verbose=1,repeat_count=1,v_smear=0.1,pos_smear=0.001,dir_smear=0.01)
\end{lstlisting}

\subsection*{Input parameters}
Parameters in \textbf{boldface} are required; the others are optional.

\begin{longtable}{p{0.22\textwidth}p{0.12\textwidth}p{0.46\textwidth}p{0.14\textwidth}}
\toprule
\textbf{Name} & \textbf{Unit} & \textbf{Description} & \textbf{Default} \\
\midrule
\endhead
filename & str & Name of neutron mcpl file to read & 0 \\
polarisationuse &   & If !=0 read polarisation vectors from file & 1 \\
verbose &   & Print debugging information for first 10 particles read & 1 \\
Emin & meV & Lower energy bound. Particles found in the MCPL-file below the limit are skipped & 0 \\
Emax & meV & Upper energy bound. Particles found in the MCPL-file above the limit are skipped & FLT\_MAX \\
repeat\_count & 1 & Repeat contents of the MCPL file this number of times. NB: When running MPI, repeating is implicit and is taken into account by integer division. MUST be combined with \_smear options! & 1 \\
v\_smear & 1 & When repeating events, make a Gaussian MC choice within v\_smear*V around particle velocity V & 0 \\
pos\_smear & m & When repeating events, make a flat MC choice of position within pos\_smear around particle starting position & 0 \\
dir\_smear & deg & When repeating events, make a Gaussian MC choice of direction within dir\_smear around particle direction & 0 \\
preload &   & Load particles during INITIALIZE. On GPU preload is forced & 0 \\
\bottomrule
\end{longtable}

\subsection*{Links}
\begin{itemize}
  \item \href{run:/home/willend/willend-McCode/mcstas-comps/misc/MCPL_input.comp}{Source code} for \texttt{MCPL\_input.comp}.
\end{itemize}
\IfFileExists{MCPL_input_static.tex}{\input{MCPL_input_static.tex}}{}

\input{misc/MCPL_output}

\newpage
\section{Progress\_bar: Simulation progress and automatic saving}
\component{Progress\_bar}{System}{percent, flag\_save, profile}{}{}
\label{s:progress-bar}
\index{Simulation progress bar}

This component displays the simulation progress and status
but does not affect the neutron parameters.
The display is updated in regular intervals of the full simulation;
the default step size is 10 \%, but it may be changed using
the \verb+percent+ parameter (from 0 to 100).
The estimated computation time is displayed at the begining
and actual simulation time is shown at the end.

Additionally, setting the \verb+flag_save+ to 1 results in
a regular save of the data files during the simulation.
This means that is is possible to view the data before the end
of the computation, and have also a trace of it in case of
computer crash. The achieved percentage of the simulation is stored in these temporary
data files. Technically, this save is equivalent to sending regularly
a USR2 signal to the running simulation.

The optional 'profile' parameter, when set to a file name, will produce the number of statistical events reaching each component in the simulation. This may be used to identify positions where events are lost.

\section{Beam\_spy: A beam analyzer}
\component{Beam\_spy}{System}{}{}{should overlap previous component}
\index{Monitors!Beam analyzer}

\textbf{Note:} \texttt{Beam\_spy} has moved to the \textbf{obsolete} component
category (see the Component Manual's Obsolete chapter); the general-purpose
\texttt{Monitor\_nD} component (Monitors chapter) now covers this
functionality and more. The description below is kept for reference.

This component is at the same time an Arm and a simple Monitor. It analyzes all neutrons reaching it, and computes statistics for the beam, as well as the intensity.

This component does not affect the neutron beam, and does not contain any propagation call. Thus it gets neutrons from the previous component in the instrument description, and should better be placed at the same position, with \verb+AT (0,0,0) RELATIVE PREVIOUS+.
